% ============================================================
% CENTRO DE CONSOLIDACIÓN LOGÍSTICA — GRUPO 20
% Presentación Junta Directiva · Marzo 2026
% Escenario ácido: $0 descuento mula G20
% Sin peajes en estos trayectos
% ============================================================
\documentclass[aspectratio=169, 10pt]{beamer}

\usepackage[utf8]{inputenc}
\usepackage[T1]{fontenc}
\usepackage[spanish]{babel}
\usepackage{booktabs}
\usepackage{colortbl}
\usepackage{tikz}
\usepackage{fontawesome5}
\usepackage{array}
\usepackage{hyperref}

\usetikzlibrary{calc, positioning, arrows.meta, shapes.geometric, fit}

\definecolor{verde}{HTML}{1A8C4E}
\definecolor{verdeclaro}{HTML}{E8F5E9}
\definecolor{rojo}{HTML}{C62828}
\definecolor{rojoclaro}{HTML}{FFEBEE}
\definecolor{azul}{HTML}{1565C0}
\definecolor{azulclaro}{HTML}{E3F2FD}
\definecolor{amarillo}{HTML}{E65100}
\definecolor{amarilloclaro}{HTML}{FFF3E0}
\definecolor{gris}{HTML}{546E7A}
\definecolor{grisclaro}{HTML}{ECEFF1}
\definecolor{negro}{HTML}{212121}
\definecolor{naranja}{HTML}{FF7A30}
\definecolor{cian}{HTML}{0097A7}

\usetheme{Madrid}
\usecolortheme{default}
\setbeamercolor{structure}{fg=verde}
\setbeamercolor{frametitle}{bg=verde!8, fg=negro}
\setbeamercolor{footline}{bg=grisclaro, fg=gris}
\setbeamercolor{block title}{bg=verde, fg=white}
\setbeamercolor{block body}{bg=verdeclaro}
\setbeamercolor{block title alerted}{bg=rojo, fg=white}
\setbeamercolor{block body alerted}{bg=rojoclaro}
\setbeamercolor{block title example}{bg=azul, fg=white}
\setbeamercolor{block body example}{bg=azulclaro}
\setbeamertemplate{navigation symbols}{}
\setbeamertemplate{itemize item}{\color{verde}\scriptsize$\blacksquare$}

\setbeamertemplate{footline}{%
    \leavevmode\hbox{%
        \begin{beamercolorbox}[wd=.5\paperwidth,ht=2.5ex,dp=1ex,left]{footline}%
            \hspace{8pt}{\tiny\color{gris} Grupo 20 $\cdot$ Junta Directiva $\cdot$ Marzo 2026}
        \end{beamercolorbox}%
        \begin{beamercolorbox}[wd=.5\paperwidth,ht=2.5ex,dp=1ex,right]{footline}%
            {\tiny\color{gris} \insertframenumber/\inserttotalframenumber}\hspace{8pt}
        \end{beamercolorbox}%
    }%
}

\newcommand{\tagline}[2]{{\scriptsize\color{#1}\textbf{\MakeUppercase{#2}}}}

\begin{document}

% ============================================================
% 1. PORTADA
% ============================================================
\begin{frame}[plain]
\vfill
\begin{center}
\begin{tikzpicture}
    \node[rounded corners=5pt, fill=verde, minimum width=0.9cm, minimum height=0.9cm,
          font=\small\bfseries, text=white] {G20};
\end{tikzpicture}
\vspace{4pt}

{\Large\color{verde}\textbf{ESTUDIO DE VIABILIDAD}}
\vspace{2pt}

{\fontsize{24}{28}\selectfont\textbf{Centro de Consolidaci\'on Log\'istica}\par}
\vspace{2pt}
\rule{3cm}{1pt}
\vspace{3pt}

{\small\color{gris} Modelo 3 Nodos $\cdot$ 19 fincas $\cdot$ Escenario \'acido}\\[2pt]
{\small\color{gris} Urab\'a, Antioquia $\cdot$ Junta Directiva $\cdot$ Marzo 2026}
\vspace{10pt}

\colorbox{grisclaro}{\small\strut\hspace{4pt}Gemelo Digital: \textcolor{verde}{\textbf{gemelo-digital-uraba.vercel.app}}\hspace{4pt}}
\vspace{6pt}

{\tiny\color{gris} Datos: Anderson Mestra Mar 2026 $\cdot$ Geocercas GPS $\cdot$ OSRM routing $\cdot$ DANE/UPRA}
\end{center}
\vfill
\end{frame}

% ============================================================
% 1b. FUENTES DE DATOS (PEDAGOGICO)
% ============================================================
\begin{frame}
\frametitle{4 fuentes. Cada una responde una pregunta distinta.}
\tagline{azul}{Fuentes de Datos}
\vspace{-2pt}

{\scriptsize Antes de cualquier n\'umero, as\'i es como est\'a construida la evidencia. Ninguna cifra del modelo sale del aire.}

\vspace{4pt}
\begin{columns}[T]
\begin{column}{0.48\textwidth}

\colorbox{azulclaro}{\parbox{\dimexpr\linewidth-8pt}{\tiny\strut
\textcolor{azul}{\textbf{1. VIAJES 2025 (Excel Anderson)}}\\[2pt]
{\scriptsize\textit{Viajes 2025 GrupoVeinte.xlsx}}\\
24,888 registros, 12,860 viajes \'unicos, 303 d\'ias.\\[3pt]
$\rightarrow$ \textbf{Vol\'umenes reales:} 5.06M cajas, 19 fincas, 485 placas, 691 conductores\\
$\rightarrow$ \textbf{Ocupaci\'on real:} 393 cajas/viaje = 72.8\% (no 76\%)\\
$\rightarrow$ \textbf{Patr\'on horario:} 39.6\% fuera de hora, 2.3\% \'optimas
\strut}}

\vspace{6pt}
\colorbox{rojoclaro}{\parbox{\dimexpr\linewidth-8pt}{\tiny\strut
\textcolor{rojo}{\textbf{2. GEOCERCAS GPS (10 meses)}}\\[2pt]
{\scriptsize\textit{geocercasvehiculos.zip}}\\
132,488 eventos GPS, 16 veh\'iculos, Mar--Dic 2025.\\[3pt]
$\rightarrow$ \textbf{Tiempo en finca:} mediana 14 min, media 53 min, P90 1h37m\\
$\rightarrow$ \textbf{Tiempo en Zungo:} mediana 1h53m, P90 4h09m\\
$\rightarrow$ \textbf{18\% de viajes} con espera $>$ 1 hora (2,315 viajes/a\~no)
\strut}}

\end{column}
\begin{column}{0.48\textwidth}

\colorbox{amarilloclaro}{\parbox{\dimexpr\linewidth-8pt}{\tiny\strut
\textcolor{amarillo}{\textbf{3. REUNIONES OP. (Marzo 2026)}}\\[2pt]
{\scriptsize\textit{Anderson Mestra + Juan Pablo L\'opez}}\\
Talleres 31 Mar 2026 (2 transcripciones).\\[3pt]
$\rightarrow$ \textbf{Tarifas reales:} sencillo \$300K Chinita (10p), mula \$720K (24p)\\
$\rightarrow$ \textbf{Tarifa standby:} \$20K--\$30K/hora (cooperativas)\\
$\rightarrow$ \textbf{Pallets nocturnos:} $\sim$50K/a\~no $\times$ \$4,200 = \$210M/a\~no
\strut}}

\vspace{6pt}
\colorbox{verdeclaro}{\parbox{\dimexpr\linewidth-8pt}{\tiny\strut
\textcolor{cian}{\textbf{4. INFRAESTRUCTURA (Arq.)}}\\[2pt]
{\scriptsize\textit{INPUTS CENTRO DE CONSOLIDACI\'ON.pdf}}\\
Evaluaci\'on profesional, precios 2026.\\[3pt]
$\rightarrow$ \textbf{\'Area construida:} 1,800 m$^2$ (bodega 1,000 + muelles 250 + etc)\\
$\rightarrow$ \textbf{Costo m$^2$:} \$1.6--\$2.2M (bodega), \$2.2--\$3.0M (admin)\\
$\rightarrow$ \textbf{CAPEX total:} \$4,020M (vs \$1,200M estimado inicial)
\strut}}

\end{column}
\end{columns}

\vspace{4pt}
\colorbox{azulclaro}{\parbox{\dimexpr\textwidth-12pt}{\tiny\strut
\textcolor{azul}{\textbf{Trazabilidad total:}} Cada n\'umero en esta presentaci\'on se puede rastrear a una de estas 4 fuentes.
Donde hay supuesto, est\'a marcado expl\'icitamente.\strut}}
\end{frame}

% ============================================================
% 2. EL PROBLEMA
% ============================================================
\begin{frame}
\frametitle{19 fincas, cada una por su cuenta}
\tagline{rojo}{El Problema}
\vspace{-2pt}

{\scriptsize Cada finca despacha de forma independiente. Camiones a medio llenar, en una v\'ia de un solo carril.}

\vspace{4pt}
\begin{columns}[T]
\begin{column}{0.55\textwidth}
\begin{center}
\scriptsize
\renewcommand{\arraystretch}{1.0}
\begin{tabular}{>{\color{rojo}\bfseries}r l}
    \toprule
    42.4 & viajes/d\'ia en v\'ia \\
    72.8\% & llenado promedio (real 2025) \\
    5--6h & espera en puerto por mula \\
    \$406 & \$/caja promedio hoy \\
    39.6\% & llegadas fuera de hora \\
    \bottomrule
\end{tabular}
\end{center}

\vspace{4pt}
\begin{alertblock}{\scriptsize Lo que pasa hoy}
\scriptsize
\begin{itemize}\setlength\itemsep{-1pt}
    \item[\textcolor{rojo}{\scriptsize$\blacksquare$}] 19 fincas despachan individualmente
    \item[\textcolor{rojo}{\scriptsize$\blacksquare$}] 26 viajes solo de las 12 fincas Centro
    \item[\textcolor{rojo}{\scriptsize$\blacksquare$}] Flota: 4 mulas + 6 sencillos
    \item[\textcolor{rojo}{\scriptsize$\blacksquare$}] No hay peajes en estos trayectos
\end{itemize}
\end{alertblock}
\end{column}
\begin{column}{0.42\textwidth}
\colorbox{grisclaro}{\parbox{\dimexpr\linewidth-8pt}{\scriptsize\strut
\textcolor{gris}{\textbf{Fuente:}}\\
Geocercas GPS, 7 veh\'iculos\\
9--12 Marzo 2026 (IP Uniban)\\
OSRM routing real (39 rutas)\\
Anderson Mestra, Marzo 2026\strut}}
\end{column}
\end{columns}
\end{frame}

% ============================================================
% 3. LA AMENAZA TARIFARIA ($/caja y $/viaje)
% ============================================================
\begin{frame}
\frametitle{El modelo tarifario va a cambiar}
\tagline{amarillo}{La Amenaza}
\vspace{-2pt}

{\scriptsize Los transportadores exigen pasar de tarifa plana (\$/caja) a cobro por viaje. Equivalencias:}

\vspace{4pt}
\begin{center}
\scriptsize
\renewcommand{\arraystretch}{1.15}
\begin{tabular}{l c c c}
    \toprule
    & \textbf{Hoy (\$/caja)} & \textbf{Nuevo (\$/viaje)} & \textbf{Cambio} \\
    \midrule
    \textbf{Por caja} & \textcolor{amarillo}{\textbf{\$406}} & \textcolor{rojo}{\textbf{\$919}} & \textcolor{rojo}{+126\%} \\[2pt]
    \textbf{Por viaje equiv.} & \textcolor{amarillo}{\textbf{\$200K}} & \textcolor{rojo}{\textbf{\$474K}} & \textcolor{rojo}{+137\%} \\
    \bottomrule
\end{tabular}
\end{center}

\vspace{2pt}
{\tiny\color{gris} \$/caja: promedio ponderado 12 fincas Centro. \$/viaje: sencillo 12p $\times$ 54 cajas = 648 cajas/viaje.}

\vspace{4pt}
\begin{center}
\scriptsize
\renewcommand{\arraystretch}{0.9}
\begin{tabular}{l r r r r}
    \toprule
    \textbf{Ejemplo fincas} & \textbf{Hoy \$/cj} & \textbf{Hoy \$/vj} & \textbf{PA \$/vj} & \textbf{PA \$/cj} \\
    \midrule
    Chinita (17 km PA)       & \$499 & \$252K & \$360K & \$712 \\
    Maritu (30 km PA)        & \$453 & \$229K & \$444K & \$921 \\
    Maria Bonita (36 km PA)  & \$304 & \$176K & \$540K & \$1,003 \\
    \bottomrule
\end{tabular}
\end{center}

\vspace{2pt}
\colorbox{amarilloclaro}{\parbox{\dimexpr\textwidth-12pt}{\scriptsize\strut
\textcolor{amarillo}{\textbf{Los transportadores llevan 2 a\~nos exigiendo este cambio.}}
El SICE-TAC del MinTransporte ya lo regula.
La pregunta no es \textit{si} va a cambiar, sino \textit{cu\'ando}.\strut}}
\end{frame}

% ============================================================
% 4. TARIFA SENCILLO: PASO A PASO
% ============================================================
\begin{frame}
\frametitle{Paso a paso: tarifa de un cami\'on sencillo}
\tagline{azul}{Entendiendo las Tarifas (1/2)}
\vspace{-2pt}

{\scriptsize La tabla de Anderson Mestra da el precio por viaje para \textbf{10 pallets}. El sencillo carga \textbf{12 pallets}.}

\vspace{6pt}
\begin{center}
\begin{tikzpicture}[
    paso/.style={rounded corners=4pt, draw=#1!50, fill=#1!8,
        minimum width=10cm, minimum height=0.9cm, font=\scriptsize,
        inner xsep=8pt, align=left}
]
    \node[paso=azul] (p1) at (0,0) {
        \textcolor{azul}{\textbf{A.}} Tarifa tabla Anderson (10 pallets)
        \hfill \textbf{var\'ia seg\'un finca y distancia}
    };
    \node[font=\tiny, text=gris, right=2pt of p1.east, anchor=west] {Ej: \$300K};

    \node[paso=azul, below=6pt of p1] (p2) {
        \textcolor{azul}{\textbf{B.}} Sencillo carga 12 pallets, no 10.
        Factor: 12 $\div$ 10 = \textbf{$\times$1.2}
        \hfill tarifa base $\times$ 1.2
    };
    \node[font=\tiny, text=gris, right=2pt of p2.east, anchor=west] {\$300K $\times$ 1.2 = \textbf{\$360K}};

    \node[paso=verde, below=6pt of p2] (p3) {
        \textcolor{verde}{\textbf{C.}} Cada viaje carga 12 pallets $\times$ 54 cajas = \textbf{648 cajas}
        \hfill \$/caja = \$/viaje $\div$ 648
    };
    \node[font=\tiny, text=gris, right=2pt of p3.east, anchor=west] {\$360K $\div$ 648 = \textbf{\$556/caja}};
\end{tikzpicture}
\end{center}

\vspace{6pt}
\begin{center}
\scriptsize
\renewcommand{\arraystretch}{0.95}
\begin{tabular}{l r r r r}
    \toprule
    \textbf{Finca} & \textbf{Tabla (10p)} & \textbf{$\times$1.2 (12p)} & \textbf{\$/caja} & \textbf{km a PA} \\
    \midrule
    Oasis (Mini-Hub)      & \$250K & \$300K & \$437 & 9 km \\
    Chinita (Centro)      & \$300K & \$360K & \$556 & 17 km \\
    Sierra (Centro)       & \$370K & \$444K & \$685 & 25 km \\
    Tagual (Centro)       & \$400K & \$480K & \$741 & 25 km \\
    M.\ Bonita (Centro)   & \$450K & \$540K & \$833 & 36 km \\
    Pto.\ Fijo (Zungo)    & \$500K & \$600K & \$926 & 52 km \\
    \bottomrule
\end{tabular}
\end{center}
\vspace{-2pt}
{\tiny\color{gris} Fuente: tabla Anderson Mestra, Marzo 2026. Pallet: 54 cajas. No hay peajes en estos trayectos.}
\end{frame}

% ============================================================
% 5. TARIFA MULA: PASO A PASO
% ============================================================
\begin{frame}
\frametitle{Paso a paso: tarifa de una mula (tractomula)}
\tagline{azul}{Entendiendo las Tarifas (2/2)}
\vspace{-2pt}

{\scriptsize La mula carga el \textbf{doble} de pallets. La tarifa es \textbf{proporcional} al sencillo.}

\vspace{6pt}
\begin{center}
\begin{tikzpicture}[
    paso/.style={rounded corners=4pt, draw=#1!50, fill=#1!8,
        minimum width=10cm, minimum height=0.9cm, font=\scriptsize,
        inner xsep=8pt, align=left}
]
    \node[paso=naranja] (p1) at (0,0) {
        \textcolor{naranja}{\textbf{A.}} Tarifa mula (20 pallets) = tarifa sencillo (10p) $\times$ 2
        \hfill proporcional
    };
    \node[font=\tiny, text=gris, right=2pt of p1.east, anchor=west] {Ej: \$300K $\times$ 2 = \textbf{\$600K}};

    \node[paso=naranja, below=6pt of p1] (p2) {
        \textcolor{naranja}{\textbf{B.}} Mula carga 24 pallets, no 20.
        Factor: 24 $\div$ 20 = \textbf{$\times$1.2}
        \hfill tarifa mula $\times$ 1.2
    };
    \node[font=\tiny, text=gris, right=2pt of p2.east, anchor=west] {\$600K $\times$ 1.2 = \textbf{\$720K}};

    \node[paso=verde, below=6pt of p2] (p3) {
        \textcolor{verde}{\textbf{C.}} Cada mula carga 24 pallets $\times$ 54 cajas = \textbf{1,296 cajas}
        \hfill \$/caja = \$/viaje $\div$ 1,296
    };
    \node[font=\tiny, text=gris, right=2pt of p3.east, anchor=west] {\$720K $\div$ 1,296 = \textbf{\$556/caja}};

    \node[paso=verde, below=6pt of p3] (p4) {
        \textcolor{verde}{\textbf{D.}} G20 puede negociar descuento: hasta \$150K menos por viaje
        \hfill \'acido: \$0
    };
    \node[font=\tiny, text=gris, right=2pt of p4.east, anchor=west] {\$720K $-$ \$0 = \textbf{\$720K}};
\end{tikzpicture}
\end{center}

\vspace{4pt}
\colorbox{amarilloclaro}{\parbox{\dimexpr\textwidth-12pt}{\scriptsize\strut
\textcolor{amarillo}{\textbf{Supuesto clave:}} ``La tarifa es proporcional entre sencillo y mula, con tolerancia de hasta \$150K
menos, seg\'un la capacidad de negociaci\'on de G20.'' --- Informaci\'on del cliente, Marzo 2026.\strut}}

\vspace{2pt}
{\tiny\color{gris} Sencillo: max 12p (tarifa base 10p). Mula: max 24p (tarifa base 20p). Pallet: 54 cajas. Sin peajes.}
\end{frame}

% ============================================================
% 6. EJEMPLO CHINITA: SENCILLO → MULA → HUB
% ============================================================
\begin{frame}
\frametitle{Ejemplo completo: Chinita $\rightarrow$ PA}
\tagline{azul}{De la Tabla Anderson al Costo por Caja}
\vspace{-2pt}

{\scriptsize Chinita es donde se ubicar\'ia el Hub. Distancia a PA: 17.5 km. Tarifa Anderson: \$300K (10 pallets).}

\vspace{4pt}
\begin{columns}[T]
\begin{column}{0.47\textwidth}
{\scriptsize\color{amarillo}\textbf{SENCILLO (sin hub)}}
\vspace{2pt}

\scriptsize
\renewcommand{\arraystretch}{1.05}
\begin{tabular}{l r}
    \toprule
    Tarifa Anderson (10p) & \$300K \\
    $\times$ 1.2 (12 pallets) & \$360K \\
    \midrule
    \textbf{\$/viaje} & \textbf{\$360K} \\
    $\div$ 648 cajas & \textbf{\$556/caja} \\
    \bottomrule
\end{tabular}
\end{column}

\begin{column}{0.47\textwidth}
{\scriptsize\color{naranja}\textbf{MULA (con hub, \'acido \$0 desc.)}}
\vspace{2pt}

\scriptsize
\renewcommand{\arraystretch}{1.05}
\begin{tabular}{l r}
    \toprule
    Tarifa sencillo $\times$ 2 (20p) & \$600K \\
    $\times$ 1.2 (24 pallets) & \$720K \\
    $-$ descuento G20 (\'acido) & \$0 \\
    \midrule
    \textbf{\$/viaje mula} & \textbf{\$720K} \\
    $\div$ 1,296 cajas & \textbf{\$556/caja} \\
    \bottomrule
\end{tabular}
\end{column}
\end{columns}

\vspace{6pt}
\colorbox{verdeclaro}{\parbox{\dimexpr\textwidth-12pt}{\scriptsize\strut
\textcolor{verde}{\textbf{Observaci\'on:}} A \$0 de descuento, el costo \textbf{por caja} es id\'entico (\$556)
entre sencillo y mula. La ventaja de la mula es que \textbf{carga el doble}:
menos viajes, menos tr\'afico, mejor llenado. Con descuento, la mula
es m\'as barata por caja.\strut}}

\vspace{2pt}
{\tiny\color{gris} Nota: \$556/caja es solo el flete Chinita$\rightarrow$PA. El costo total del hub incluye consolidaci\'on y finca$\rightarrow$hub.}
\end{frame}

% ============================================================
% 7. LOS 3 GRUPOS
% ============================================================
\begin{frame}
\frametitle{Las 19 fincas se dividen en 3 grupos}
\tagline{azul}{Tres Destinos}
\vspace{-2pt}

\begin{center}
\scriptsize
\renewcommand{\arraystretch}{1.1}
\begin{tabular}{l c r l l}
    \toprule
    \textbf{Grupo} & \textbf{Fincas} & \textbf{Cajas/a\~no} & \textbf{Destino} & \textbf{Hub?} \\
    \midrule
    \textcolor{verde}{\textbf{Centro}} & 12 & 3.30M (63\%) & Hub $\rightarrow$ PA & \textcolor{verde}{\textbf{S\'I}} \\
    \textcolor{cian}{\textbf{Mini-Hub}} & 3 & 0.46M (9\%) & PA directo & No \\
    \textcolor{naranja}{\textbf{Zungo}} & 4 & 1.50M (28\%) & Embarcadero & No \\
    \midrule
    \textbf{Total} & \textbf{19} & \textbf{5.26M} & & \\
    \bottomrule
\end{tabular}
\end{center}

\vspace{4pt}
\begin{columns}[T]
\begin{column}{0.30\textwidth}
\colorbox{verdeclaro}{\parbox{\dimexpr\linewidth-8pt}{\scriptsize\strut
\textcolor{verde}{\textbf{12 Centro:}} Lejos de PA.
26 viajes individuales/d\'ia.
\textbf{El hub consolida en 10 mulas.}\strut}}
\end{column}
\begin{column}{0.30\textwidth}
\colorbox{azulclaro}{\parbox{\dimexpr\linewidth-8pt}{\scriptsize\strut
\textcolor{azul}{\textbf{3 Mini-Hub:}} Oasis, Dunas,
Guadalupe. \textbf{Cerca de PA.}
Despacho directo.\strut}}
\end{column}
\begin{column}{0.30\textwidth}
\colorbox{amarilloclaro}{\parbox{\dimexpr\linewidth-8pt}{\scriptsize\strut
\textcolor{amarillo}{\textbf{4 Zungo:}} Pto.\ Fijo,
Almendros, Matogroso 1 y 2.
\textbf{Exportan v\'ia fluvial.}\strut}}
\end{column}
\end{columns}

\vspace{3pt}
{\tiny\color{gris} El ahorro del hub solo aplica a las 12 fincas Centro (3.30M cajas/a\~no, 63\% de G20).}
\end{frame}

% ============================================================
% 7b. LAS 19 FINCAS (DETALLE)
% ============================================================
\begin{frame}
\frametitle{Las 19 fincas: producci\'on y costos}
\tagline{azul}{Detalle por Finca}
\vspace{-4pt}

\begin{center}
\tiny
\renewcommand{\arraystretch}{0.88}
\begin{tabular}{c l l r r r r}
    \toprule
    \textbf{ID} & \textbf{Finca} & \textbf{Grupo} & \textbf{Cajas/sem} & \textbf{Pall/d} & \textbf{\$/cj hoy} & \textbf{\$/vj PA} \\
    \midrule
    \rowcolor{amarilloclaro}
    F01 & Punto Fijo    & \textcolor{naranja}{Zungo}    & 7,856 & 29.1 & \$359 & \$500K \\
    \rowcolor{amarilloclaro}
    F02 & Almendros     & \textcolor{naranja}{Zungo}    & 5,931 & 22.0 & \$349 & \$500K \\
    \rowcolor{azulclaro}
    F03 & Oasis         & \textcolor{cian}{Mini-Hub}     & 3,430 & 12.7 & \$519 & \$250K \\
    \rowcolor{azulclaro}
    F04 & Dunas         & \textcolor{cian}{Mini-Hub}     & 2,711 & 10.0 & \$519 & \$250K \\
    \rowcolor{azulclaro}
    F05 & Guadalupe     & \textcolor{cian}{Mini-Hub}     & 2,667 &  9.9 & \$519 & \$250K \\
    \rowcolor{verdeclaro}
    F06 & Chinita       & \textcolor{verde}{Centro}      & 10,113 & 37.5 & \$499 & \$300K \\
    \rowcolor{verdeclaro}
    F07 & Santa Marta   & \textcolor{verde}{Centro}      & 7,731 & 28.6 & \$499 & \$300K \\
    \rowcolor{verdeclaro}
    F08 & La Llave      & \textcolor{verde}{Centro}      & 4,104 & 15.2 & \$499 & \$300K \\
    \rowcolor{amarilloclaro}
    F09 & Matogroso 2   & \textcolor{naranja}{Zungo}     & 8,375 & 31.0 & \$295 & \$500K \\
    \rowcolor{amarilloclaro}
    F10 & Matogroso 1   & \textcolor{naranja}{Zungo}     & 6,662 & 24.7 & \$295 & \$500K \\
    \rowcolor{verdeclaro}
    F11 & Sierra        & \textcolor{verde}{Centro}      & 2,784 & 10.3 & \$453 & \$370K \\
    \rowcolor{verdeclaro}
    F12 & Villa Fresia  & \textcolor{verde}{Centro}      & 3,071 & 11.4 & \$453 & \$370K \\
    \rowcolor{verdeclaro}
    F13 & Tagual        & \textcolor{verde}{Centro}      & 4,464 & 16.5 & \$453 & \$400K \\
    \rowcolor{verdeclaro}
    F14 & San Felipe    & \textcolor{verde}{Centro}      & 4,982 & 18.5 & \$453 & \$400K \\
    \rowcolor{verdeclaro}
    F15 & Maritu        & \textcolor{verde}{Centro}      & 2,530 &  9.4 & \$453 & \$370K \\
    \rowcolor{verdeclaro}
    F16 & San Juan      & \textcolor{verde}{Centro}      & 4,235 & 15.7 & \$426 & \$400K \\
    \rowcolor{verdeclaro}
    F17 & Santa Helena  & \textcolor{verde}{Centro}      & 3,357 & 12.4 & \$426 & \$400K \\
    \rowcolor{verdeclaro}
    F18 & Maria Bonita  & \textcolor{verde}{Centro}      & 7,827 & 29.0 & \$304 & \$450K \\
    \rowcolor{verdeclaro}
    F19 & Tucanes       & \textcolor{verde}{Centro}      & 8,173 & 30.3 & \$304 & \$450K \\
    \midrule
    & \textbf{TOTAL} & \textbf{19 fincas} & \textbf{101,003} & \textbf{374.2} & & \\
    \bottomrule
\end{tabular}
\end{center}

\vspace{1pt}
{\tiny\color{gris} \$/cj hoy = tarifa a Zungo (tz). \$/vj PA = tarifa Anderson 10 pallets. Pall/d = cajas/sem $\div$ 5 $\div$ 54.
\textcolor{verde}{Verde}: Centro (12, van al hub). \textcolor{cian}{Cian}: Mini-Hub (3, PA directo). \textcolor{naranja}{Naranja}: Zungo (4, embarcadero).}
\end{frame}

% ============================================================
% 8. LA SOLUCIÓN: 3 NODOS
% ============================================================
\begin{frame}
\frametitle{Consolidar antes de despachar}
\tagline{verde}{La Soluci\'on}
\vspace{-2pt}

{\scriptsize Un Centro en \textbf{Chinita} donde la fruta de 12 fincas se paletiza, conteneriza, y sale en \textbf{mulas llenas} a PA.}

\vspace{4pt}
\begin{center}
\begin{tikzpicture}[
    nodo/.style={circle, draw=#1, fill=#1!15, minimum size=0.85cm, font=\scriptsize},
    linea/.style={-{Stealth[length=4pt]}, thick, color=#1},
    etiq/.style={font=\tiny, text=#1, align=center}
]
    \node[nodo=naranja] (f) at (0,0) {\faIcon{seedling}};
    \node[etiq=naranja, below=1pt of f] {19 Fincas};

    \node[nodo=verde] (h) at (2.8,0) {\faIcon{warehouse}};
    \node[etiq=verde, below=1pt of h] {Centro\\12 fincas};

    \node[nodo=cian] (m) at (5.5,1.2) {\faIcon{truck}};
    \node[etiq=cian, below=1pt of m] {Mini-Hub (3)};

    \node[nodo=naranja] (z) at (5.5,-1.2) {\faIcon{ship}};
    \node[etiq=naranja, below=1pt of z] {Zungo (4)};

    \node[nodo=azul] (p) at (9,0) {\faIcon{anchor}};
    \node[etiq=azul, below=1pt of p] {Puerto Ant.};

    \draw[linea=verde] (f) -- node[above, font=\tiny, text=gris] {26 senc.} (h);
    \draw[linea=verde, line width=2pt] (h) -- node[below, font=\tiny\bfseries, text=verde] {10 mulas} (p);
    \draw[linea=cian] (m) -- node[above, font=\tiny, text=gris, sloped] {4 vj} (p);
    \draw[linea=naranja] (f.south east) -- (z.west);
\end{tikzpicture}
\end{center}

\vspace{4pt}
\colorbox{verdeclaro}{\parbox{\dimexpr\textwidth-12pt}{\scriptsize\strut
\textcolor{verde}{\textbf{En la v\'ia principal:}} 41 viajes $\rightarrow$ 25 viajes (\textbf{$-$39\%}).\\
Se agregan 26 viajes cortos finca$\rightarrow$hub (8--15 km, v\'ias locales).
Contenerizado, 89--98\% llenado, espera $<$1h.\strut}}
\end{frame}

% ============================================================
% 8b. LA CUENTA REAL DE VIAJES
% ============================================================
\begin{frame}
\frametitle{La cuenta real de viajes}
\tagline{amarillo}{Transparencia: los viajes totales SUBEN}
\vspace{-2pt}

{\scriptsize El hub reduce viajes en la v\'ia principal, pero \textbf{agrega} viajes cortos finca$\rightarrow$hub:}

\vspace{4pt}
\begin{center}
\scriptsize
\renewcommand{\arraystretch}{1.15}
\begin{tabular}{l c c l}
    \toprule
    \textbf{Tipo de viaje} & \textbf{SC1 (Hoy)} & \textbf{SC3 (Hub)} & \textbf{Cambio} \\
    \midrule
    \multicolumn{4}{l}{\textcolor{rojo}{\textbf{V\'IA PRINCIPAL (carretera de un carril)}}} \\
    \midrule
    Centro $\rightarrow$ Zungo/PA & \textcolor{rojo}{26} & --- & reemplazados \\
    Hub $\rightarrow$ PA (mulas) & --- & \textcolor{verde}{10} & {\tiny nuevos, consolidados} \\
    Mini-Hub $\rightarrow$ PA & 4 & 4 & sin cambio \\
    Zungo $\rightarrow$ Emb. & 11 & 11 & sin cambio \\
    \cmidrule{2-3}
    \textbf{Subtotal v\'ia ppal.} & \textcolor{rojo}{\textbf{41}} & \textcolor{verde}{\textbf{25}} & \textcolor{verde}{\textbf{$-$39\%}} \\
    \midrule
    \multicolumn{4}{l}{\textcolor{azul}{\textbf{V\'IAS LOCALES (internas, 8--15 km)}}} \\
    \midrule
    Finca $\rightarrow$ Hub & 0 & \textcolor{azul}{26} & {\tiny nuevos, cortos} \\
    \midrule
    \textbf{TOTAL viajes} & \textbf{41} & \textbf{51} & \textcolor{rojo}{\textbf{+24\%}} \\
    \bottomrule
\end{tabular}
\end{center}

\vspace{4pt}
\colorbox{amarilloclaro}{\parbox{\dimexpr\textwidth-12pt}{\scriptsize\strut
\textcolor{amarillo}{\textbf{Lectura honesta:}} Los viajes totales \textbf{suben} de 41 a 51.
Pero los 26 nuevos viajes finca$\rightarrow$hub son \textbf{cortos} (8--15 km) y en \textbf{v\'ias locales/internas},
no en la carretera congestionada de un solo carril. El beneficio real es la
\textbf{reducci\'on de 41 a 25 en la v\'ia principal ($-$39\%)}, con carga contenerizada.\strut}}
\end{frame}

% ============================================================
% 9. LÓGICA DE LA CONSOLIDACIÓN
% ============================================================
\begin{frame}
\frametitle{\textquestiondown Por qu\'e 10 mulas en vez de 26 sencillos?}
\tagline{verde}{L\'ogica de la Consolidaci\'on}
\vspace{-2pt}

\begin{columns}[T]
\begin{column}{0.48\textwidth}
{\scriptsize\color{rojo}\textbf{SIN HUB: 26 sencillos (12p c/u)}}
\vspace{1pt}

\scriptsize
\renewcommand{\arraystretch}{0.8}
\begin{tabular}{l r r r}
    \toprule
    \textbf{Finca} & \textbf{Pall/d} & \textbf{Vj} & \textbf{Vac\'ios} \\
    \midrule
    Chinita       & 37.5 & 4 & 10.5 \\
    Tucanes       & 30.3 & 3 & 5.7 \\
    M.\ Bonita    & 29.0 & 3 & 7.0 \\
    Sta.\ Helena  & 12.4 & 2 & \textcolor{rojo}{11.6} \\
    Maritu        &  9.4 & 1 & 2.6 \\
    {\color{gris}\tiny ... 7 m\'as} & & & \\
    \midrule
    \textbf{12 Centro} & \textbf{234.8} & \textcolor{rojo}{\textbf{26}} & \textcolor{rojo}{\textbf{77}} \\
    \bottomrule
\end{tabular}

\vspace{2pt}
{\tiny 26 vj $\times$ 12p = 312 slots. Solo 235 usados.\\
\textcolor{rojo}{\textbf{77 pallets de espacio desperdiciado.}}}
\end{column}

\begin{column}{0.48\textwidth}
{\scriptsize\color{verde}\textbf{CON HUB: 10 mulas (24p c/u)}}
\vspace{1pt}

\begin{block}{\scriptsize Se consolida y se manda en mulas}
\scriptsize
234.8 pallets $\div$ 24 = 9.78 $\rightarrow$ \textbf{10 mulas}

10 $\times$ 24p = 240 slots $\rightarrow$ \textbf{solo 5 vac\'ios}

\vspace{2pt}
\textcolor{verde}{\textbf{De 77 vac\'ios a 5.}}

Los pallets sobrantes de una finca llenan el espacio de otra.
\end{block}

\vspace{2pt}
\colorbox{verdeclaro}{\parbox{\dimexpr\linewidth-8pt}{\tiny\strut
\textbf{26 sencillos $\rightarrow$ 10 mulas}\\
16 viajes menos. Mini-Hub (4) y Zungo (11) sin cambio.\\
\textbf{Total G20: 41 $\rightarrow$ 25 ($-$39\%).}\strut}}
\end{column}
\end{columns}
\end{frame}

% ============================================================
% 10. CÓMO SE ARMA EL COSTO HUB→PA (1/2)
% ============================================================
\begin{frame}
\frametitle{Costo del hub: flete Hub $\rightarrow$ PA}
\tagline{amarillo}{Paso a Paso (1/3) --- Escenario \'Acido}
\vspace{-2pt}

{\scriptsize El componente m\'as grande del costo. Paso a paso usando Chinita como referencia:}

\vspace{6pt}
\begin{center}
\begin{tikzpicture}[
    paso/.style={rounded corners=3pt, draw=#1!50, fill=#1!8,
        minimum width=10cm, minimum height=0.75cm, font=\scriptsize,
        inner xsep=8pt, align=left}
]
    \node[paso=azul] (p1) at (0,0) {
        \textcolor{azul}{\textbf{1.}} Chinita sencillo (tabla Anderson, 10p): \$300K
    };
    \node[paso=azul, below=4pt of p1] (p2) {
        \textcolor{azul}{\textbf{2.}} Tarifa mula = sencillo $\times$ 2 (proporcional): \$300K $\times$ 2 = \textbf{\$600K} (20p)
    };
    \node[paso=naranja, below=4pt of p2] (p3) {
        \textcolor{naranja}{\textbf{3.}} Mula carga 24p, no 20. Factor 1.2: \$600K $\times$ 1.2 = \textbf{\$720K} (24p)
    };
    \node[paso=amarillo, below=4pt of p3] (p4) {
        \textcolor{amarillo}{\textbf{4.}} Descuento G20 (\'acido): \$0. Total: \textbf{\$720K/mula}
    };
    \node[paso=verde, below=8pt of p4] (p5) {
        \textcolor{verde}{\textbf{5.}} Centro produce 12,674 cajas/d\'ia $\div$ 1,296 cajas/mula = \textbf{10 mulas/d\'ia}
    };
    \node[paso=verde, below=4pt of p5, minimum height=0.9cm] (p6) {
        \textcolor{verde}{\textbf{6.}} Flete/caja = 10 mulas $\times$ \$720K $\div$ 12,674 cajas = \textbf{\$568/caja}
    };
\end{tikzpicture}
\end{center}

\vspace{1pt}
{\tiny\color{gris} 12,674 cajas/d\'ia = suma 12 fincas Centro (cajas/sem $\div$ 5). 1,296 cajas/mula = 24p $\times$ 54.}
\end{frame}

% ============================================================
% 11. CONSOLIDACIÓN + FINCA→HUB (2/3 y 3/3)
% ============================================================
\begin{frame}
\frametitle{Costo del hub: consolidaci\'on + finca $\rightarrow$ hub}
\tagline{amarillo}{Paso a Paso (2/3 y 3/3)}
\vspace{-2pt}

\begin{columns}[T]
\begin{column}{0.47\textwidth}
{\scriptsize\color{azul}\textbf{COMPONENTE 2: CONSOLIDACI\'ON}}
\vspace{2pt}

\scriptsize
\renewcommand{\arraystretch}{1.1}
\begin{tabular}{l r}
    \toprule
    Paletizar & $\sim$\$15/caja \\
    Inspeccionar & $\sim$\$12/caja \\
    Cargar en contenedor & $\sim$\$15/caja \\
    \midrule
    \textcolor{azul}{\textbf{Total consolidaci\'on}} & \textcolor{azul}{\textbf{\$42/caja}} \\
    \bottomrule
\end{tabular}

\vspace{6pt}
\colorbox{azulclaro}{\parbox{\dimexpr\linewidth-8pt}{\tiny\strut
\textcolor{azul}{\textbf{Nota:}} Si vendemos servicio a terceros,
\textbf{\$42/caja es el piso} (costo de operaci\'on).
La tarifa de servicio debe ser mayor.\strut}}
\end{column}

\begin{column}{0.47\textwidth}
{\scriptsize\color{cian}\textbf{COMPONENTE 3: FINCA $\rightarrow$ HUB}}
\vspace{2pt}

\scriptsize
\renewcommand{\arraystretch}{1.1}
\begin{tabular}{l r}
    \toprule
    Distancia promedio & 8--15 km \\
    \$/km promedio (tabla) & \$13,140/km \\
    Viaje sencillo equiv. & $\sim$\$130K \\
    $\div$ 648 cajas/viaje & \\
    \midrule
    \textcolor{cian}{\textbf{Finca $\rightarrow$ Hub}} & \textcolor{cian}{\textbf{\$200/caja}} \\
    \bottomrule
\end{tabular}

\vspace{6pt}
\colorbox{amarilloclaro}{\parbox{\dimexpr\linewidth-8pt}{\tiny\strut
\textcolor{amarillo}{\textbf{Supuesto \'acido (alto).}}
\$200/caja = \$130K/viaje para $\sim$10 km.
Rango real: \$80--\$300/caja seg\'un distancia.
\textbf{Variable ajustable en gemelo digital.}\strut}}
\end{column}
\end{columns}
\end{frame}

% ============================================================
% 12. TOTAL HUB POR CAJA
% ============================================================
\begin{frame}
\frametitle{Costo total del hub por caja}
\tagline{verde}{Sumando los 3 Componentes}
\vspace{-2pt}

{\scriptsize Escenario \'acido: \$0 descuento mula, finca$\rightarrow$hub alto (\$200), tarifa sin negociar.}

\vspace{6pt}
\begin{center}
\begin{tikzpicture}[
    bloque/.style={rounded corners=3pt, minimum width=9cm, minimum height=0.8cm,
        font=\scriptsize, inner xsep=10pt, align=left}
]
    \node[bloque, draw=amarillo!50, fill=amarillo!8] (b1) at (0,0) {
        \textcolor{amarillo}{\textbf{1.}} Hub $\rightarrow$ PA (flete mula, 10 mulas/d\'ia)
        \hfill \textcolor{amarillo}{\textbf{\$568/caja}}
    };
    \node[bloque, draw=azul!40, fill=azul!6, below=4pt of b1] (b2) {
        \textcolor{azul}{\textbf{2.}} Consolidaci\'on (paletizar + inspeccionar + cargar)
        \hfill \textcolor{azul}{\textbf{\$42/caja}}
    };
    \node[bloque, draw=cian!50, fill=cian!6, below=4pt of b2] (b3) {
        \textcolor{cian}{\textbf{3.}} Finca $\rightarrow$ Hub (transporte $\sim$8--15 km)
        \hfill \textcolor{cian}{\textbf{\$200/caja}}
    };
    \node[bloque, draw=verde!60, fill=verde!10, below=8pt of b3, minimum height=1cm] (tot) {
        {\large\textbf{TOTAL HUB}} \hspace{4pt} {\scriptsize (\$568 + \$42 + \$200)}
        \hfill {\fontsize{18}{20}\selectfont\textcolor{verde}{\textbf{\$810/caja}}}
    };
\end{tikzpicture}
\end{center}

\vspace{6pt}
\begin{columns}[T]
\begin{column}{0.47\textwidth}
\colorbox{rojoclaro}{\parbox{\dimexpr\linewidth-8pt}{\scriptsize\strut
\textcolor{rojo}{\textbf{Sin hub:}} \$919/caja\\
{\tiny (26 sencillos individuales, prom.\ Centro)}\strut}}
\end{column}
\begin{column}{0.47\textwidth}
\colorbox{verdeclaro}{\parbox{\dimexpr\linewidth-8pt}{\scriptsize\strut
\textcolor{verde}{\textbf{Con hub:}} \$810/caja\\
{\tiny (10 mulas consolidadas, \'acido \$0 desc.)}\strut}}
\end{column}
\end{columns}

\vspace{4pt}
\begin{center}
{\large\textcolor{verde}{\textbf{Ahorro: \$109/caja}} {\scriptsize (\$919 $-$ \$810)}}
\end{center}
\end{frame}

% ============================================================
% 13. ESCENARIO ÁCIDO: LA CUENTA
% ============================================================
\begin{frame}
\frametitle{Escenario \'acido: la cuenta anual}
\tagline{amarillo}{El Peor Caso Posible}
\vspace{-2pt}

\begin{center}
\scriptsize
\renewcommand{\arraystretch}{1.15}
\begin{tabular}{l r l}
    \toprule
    Ahorro/caja (\'acido) & \$109 & $=$ \$919 $-$ \$810 \\
    $\times$ 3.30M cajas Centro/a\~no & & \\
    \midrule
    \textbf{Ahorro bruto} & \textbf{\$360M/a\~no} & \\
    $-$ OPEX hub & $-$\$437M & {\tiny 6 personas + svc + mant + seg + admin} \\
    \midrule
    \textcolor{rojo}{\textbf{Neto transporte solo}} & \textcolor{rojo}{\textbf{$-$\$77M/a\~no}} & \textcolor{rojo}{\textbf{Deficit con CAPEX real}} \\
    \bottomrule
\end{tabular}
\end{center}

\vspace{6pt}
\colorbox{amarilloclaro}{\parbox{\dimexpr\textwidth-12pt}{\scriptsize\strut
\textcolor{amarillo}{\textbf{Supuestos del peor caso:}}
\begin{itemize}\setlength\itemsep{-2pt}\scriptsize
    \item[\textcolor{amarillo}{\tiny$\blacksquare$}] Descuento mula G20: \textbf{\$0} (no se logra negociar nada)
    \item[\textcolor{amarillo}{\tiny$\blacksquare$}] Finca$\rightarrow$Hub: \textbf{\$200/caja} (supuesto alto)
    \item[\textcolor{amarillo}{\tiny$\blacksquare$}] Tarifa Hub$\rightarrow$PA: \textbf{\$720K/mula} (sin negociar)
    \item[\textcolor{amarillo}{\tiny$\blacksquare$}] CAPEX real: \textbf{\$4,020M} (1,800 m$^2$ infraestructura completa)
    \item[\textcolor{amarillo}{\tiny$\blacksquare$}] OPEX: \textbf{\$437M/a\~no} (escalado a nueva infraestructura)
\end{itemize}
\strut}}

\vspace{2pt}
\colorbox{rojoclaro}{\parbox{\dimexpr\textwidth-12pt}{\scriptsize\strut
\textcolor{rojo}{\textbf{Con CAPEX real, transporte solo no alcanza.}}
Pero este escenario ignora 3 palancas cuantificadas:
\textcolor{cian}{\textbf{standby evitado}}, \textcolor{cian}{\textbf{nocturnos evitados}}, y
\textcolor{verde}{\textbf{servicio a terceros}}. Ver siguientes slides.\strut}}
\end{frame}

% ============================================================
% 13b. DATOS REALES 2025
% ============================================================
\begin{frame}
\frametitle{12,860 viajes procesados confirman el problema}
\tagline{cian}{Datos Reales 2025}
\vspace{-2pt}

\begin{columns}[T]
\begin{column}{0.48\textwidth}
\begin{center}
\scriptsize
\renewcommand{\arraystretch}{1.1}
\begin{tabular}{>{\color{cian}\bfseries}r l}
    \toprule
    12,860 & viajes / a\~no \\
    42.4 & viajes / d\'ia promedio \\
    5.06M & cajas remitidas \\
    5.35M & cajas recibidas (+5.8\%) \\
    \rowcolor{rojoclaro}
    72.8\% & ocupaci\'on real \\
    \rowcolor{rojoclaro}
    39.6\% & llegadas fuera de hora \\
    485 & veh\'iculos, 691 conductores \\
    \bottomrule
\end{tabular}
\end{center}

\vspace{4pt}
\colorbox{azulclaro}{\parbox{\dimexpr\linewidth-8pt}{\tiny\strut
\textcolor{azul}{\textbf{Fuente:}} Viajes 2025 GrupoVeinte.xlsx\\
24,888 registros, 303 d\'ias operativos, 19 fincas.\strut}}
\end{column}

\begin{column}{0.48\textwidth}
{\scriptsize\color{gris}\textbf{TOP FINCAS POR VOLUMEN}}
\vspace{2pt}

\scriptsize
\renewcommand{\arraystretch}{1.0}
\begin{tabular}{l r r}
    \toprule
    \textbf{Finca} & \textbf{Cajas} & \textbf{Ha} \\
    \midrule
    Chinita      & 525K & 235 \\
    Matogroso 2  & 410K & 213 \\
    Santa Marta  & 398K & 133 \\
    Tucanes      & 383K & 136 \\
    Punto Fijo   & 370K & 171 \\
    \bottomrule
\end{tabular}

\vspace{4pt}
{\scriptsize\color{gris}\textbf{PATRONES CR\'ITICOS}}
\begin{itemize}\setlength\itemsep{-2pt}\tiny
    \item[\textcolor{rojo}{\tiny$\blacksquare$}] 52\% salidas entre 16:00--20:00
    \item[\textcolor{rojo}{\tiny$\blacksquare$}] Tr\'ansito medio: 2.2h (mediana 1h)
    \item[\textcolor{rojo}{\tiny$\blacksquare$}] Top 5 veh\'iculos = 44\% viajes
    \item[\textcolor{rojo}{\tiny$\blacksquare$}] Martes = d\'ia m\'as cargado (23\%)
\end{itemize}
\end{column}
\end{columns}
\end{frame}

% ============================================================
% 13c. TIEMPOS MUERTOS (GEOCERCAS)
% ============================================================
\begin{frame}
\frametitle{132,488 eventos GPS cuantifican la espera}
\tagline{rojo}{Tiempos Muertos (Geocercas 2025)}
\vspace{-2pt}

\begin{columns}[T]
\begin{column}{0.32\textwidth}
\colorbox{rojoclaro}{\parbox{\dimexpr\linewidth-8pt}{\centering\strut
{\large\textcolor{rojo}{\textbf{53 min}}}\\[1pt]
{\tiny espera media en finca}\\
{\tiny mediana: 14 min, P90: 1h37m}\strut}}
\end{column}
\begin{column}{0.32\textwidth}
\colorbox{rojoclaro}{\parbox{\dimexpr\linewidth-8pt}{\centering\strut
{\large\textcolor{rojo}{\textbf{1h53m}}}\\[1pt]
{\tiny mediana en Zungo}\\
{\tiny P90: 4h09m --- cuello}\strut}}
\end{column}
\begin{column}{0.32\textwidth}
\colorbox{rojoclaro}{\parbox{\dimexpr\linewidth-8pt}{\centering\strut
{\large\textcolor{rojo}{\textbf{18\%}}}\\[1pt]
{\tiny viajes con espera $>$ 1h}\\
{\tiny $\sim$2,315 viajes/a\~no}\strut}}
\end{column}
\end{columns}

\vspace{6pt}
{\scriptsize\color{gris}\textbf{PEORES FINCAS G20 POR ESPERA MEDIANA}}
\vspace{2pt}

\begin{center}
\scriptsize
\renewcommand{\arraystretch}{1.05}
\begin{tabular}{l r r r r}
    \toprule
    \textbf{Finca} & \textbf{Visitas} & \textbf{Mediana} & \textbf{P90} & \textbf{Costo @\$20K/h} \\
    \midrule
    Nueva Colonia & 393 & \textcolor{rojo}{1h19m} & 3h54m & \textcolor{rojo}{\$16,333} \\
    Castillete    & 471 & \textcolor{rojo}{55m}   & 2h21m & \textcolor{rojo}{\$8,333} \\
    Tucanes       & 597 & \textcolor{amarillo}{47m}  & 2h08m & \textcolor{amarillo}{\$5,667} \\
    San Juan      & 482 & \textcolor{amarillo}{43m}  & 2h15m & \textcolor{amarillo}{\$4,333} \\
    Chinita       & 509 & \textcolor{amarillo}{40m}  & 1h57m & \textcolor{amarillo}{\$3,333} \\
    \bottomrule
\end{tabular}
\end{center}

\vspace{4pt}
\colorbox{rojoclaro}{\parbox{\dimexpr\textwidth-12pt}{\scriptsize\strut
\textcolor{rojo}{\textbf{Riesgo regulatorio:}} SICE-TAC regula tiempos muertos.
A \$20K--\$30K/hora, el standby total es \textbf{\$42--83M COP/a\~no}.
El embarcadero Zungo agrega \textbf{\$163M} potenciales.
Total exposici\'on: \textcolor{rojo}{\textbf{\$218--246M/a\~no}}.\strut}}
\end{frame}

% ============================================================
% 13d. CALCULO SERVILLETA 1/2: STANDBY + NOCTURNOS
% ============================================================
\begin{frame}
\frametitle{C\'omo llegamos a \$142M en palancas operativas}
\tagline{cian}{C\'alculo de Servilleta (1/2)}
\vspace{-2pt}

{\tiny Dos palancas cuantificables con datos reales. Cada cifra es multiplicaci\'on simple, sin modelos complejos.}

\vspace{4pt}
\begin{columns}[T]
\begin{column}{0.48\textwidth}
{\tiny\color{cian}\textbf{PALANCA A: STANDBY EVITADO}}
\vspace{2pt}

\tiny
\renewcommand{\arraystretch}{1.05}
\begin{tabular}{l r}
    \toprule
    Viajes con espera $>$ 1h (Geocercas: 18\% $\times$ 12,860) & \textcolor{cian}{\textbf{2,315}} \\
    $\times$ Espera promedio (exceso sobre 30 min cargue) & \textcolor{cian}{\textbf{0.92 h}} \\
    $\times$ Tarifa standby (rango cooperativas) & \textcolor{cian}{\textbf{\$20,000/h}} \\
    \midrule
    \textbf{= Exposici\'on total} (2,315 $\times$ 0.92 $\times$ \$20K) & \textcolor{cian}{\textbf{\$42.6M}} \\
    \rowcolor{verdeclaro}
    \textbf{$\times$ 50\% evitado con hub} (conservador) & \textcolor{cian}{\textbf{\$21M}} \\
    \bottomrule
\end{tabular}

\vspace{3pt}
{\tiny\color{gris} Escenario optimista (\$30K, 100\%): \textcolor{cian}{\textbf{\$128M/a\~no}}.}
\end{column}

\begin{column}{0.48\textwidth}
{\tiny\color{cian}\textbf{PALANCA B: PALLETS NOCTURNOS}}
\vspace{2pt}

\tiny
\renewcommand{\arraystretch}{1.05}
\begin{tabular}{l r}
    \toprule
    Pallets nocturnos / a\~no (400/d $\times$ 5 $\times$ 52 $\times$ 48\%) & \textcolor{cian}{\textbf{50,000}} \\
    $\times$ Costo nocturno por pallet (tarifa 2026) & \textcolor{cian}{\textbf{\$4,200}} \\
    \midrule
    \textbf{= Costo total nocturnos} (50K $\times$ \$4,200) & \textcolor{cian}{\textbf{\$210M}} \\
    \rowcolor{verdeclaro}
    \textbf{$\times$ 50\% capturado con hub} (JP: ``mitad'') & \textcolor{cian}{\textbf{\$105M}} \\
    \bottomrule
\end{tabular}

\vspace{3pt}
{\tiny\color{gris} El hub opera en turnos (hasta 22h) evitando recogida nocturna en finca.}
\end{column}
\end{columns}

\vspace{4pt}
\colorbox{verdeclaro}{\parbox{\dimexpr\textwidth-12pt}{\scriptsize\strut
\textcolor{cian}{\textbf{Suma palancas operativas:}} \$21M (standby) + \$105M (nocturnos) = \textbf{\$126M/a\~no}.
Con escenario medio (\$20K/h, 50\% captura ambos): estamos usando \textbf{\$142M/a\~no} en el modelo.
Fuente: Geocercas + reuniones Mar 2026.\strut}}
\end{frame}

% ============================================================
% 13e. CALCULO SERVILLETA 2/2: TERCEROS
% ============================================================
\begin{frame}
\frametitle{C\'omo llegamos a \$652M en servicio a terceros}
\tagline{verde}{C\'alculo de Servilleta (2/2)}
\vspace{-2pt}

{\tiny Con el hub construido, sobra capacidad para consolidar fruta de fincas vecinas (402 bananeras en zona). Ingreso adicional puro.}

\vspace{4pt}
\begin{columns}[T]
\begin{column}{0.48\textwidth}
{\tiny\color{verde}\textbf{SUPUESTOS (SERVILLETA)}}
\vspace{2pt}

\tiny
\renewcommand{\arraystretch}{1.05}
\begin{tabular}{l r}
    \toprule
    Producci\'on regional anual (Urab\'a) & \textcolor{verde}{\textbf{114M cajas}} \\
    $\times$ Penetraci\'on objetivo (conservador) & \textcolor{verde}{\textbf{10\%}} \\
    \midrule
    \textbf{= Cajas terceros / a\~no} (114M $\times$ 10\%) & \textcolor{verde}{\textbf{11.4M}} \\
    $\times$ Tarifa servicio (piso \$42 + margen) & \textcolor{verde}{\textbf{\$57/caja}} \\
    \rowcolor{verdeclaro}
    \textbf{= Ingreso terceros / a\~no} (11.4M $\times$ \$57) & \textcolor{verde}{\textbf{\$652M}} \\
    \bottomrule
\end{tabular}
\end{column}

\begin{column}{0.48\textwidth}
{\tiny\color{verde}\textbf{MATRIZ DE SENSIBILIDAD}}
\vspace{2pt}

\tiny
\renewcommand{\arraystretch}{1.05}
\begin{tabular}{l r r r}
    \toprule
    \textbf{Tarifa} & \textbf{5\%} & \textbf{10\%} & \textbf{25\%} \\
    \midrule
    \rowcolor{rojoclaro}
    {\color{gris}\$42 (costo)} & {\color{gris}\$239M} & {\color{gris}\$479M} & {\color{gris}\$1,197M} \\
    \$50 & \$285M & \$570M & \$1,425M \\
    \rowcolor{verdeclaro}
    \textcolor{verde}{\textbf{\$57 base}} & \$325M & \textcolor{verde}{\textbf{\$652M}} & \$1,625M \\
    \$80 & \$456M & \$912M & \$2,280M \\
    \bottomrule
\end{tabular}

\vspace{4pt}
\colorbox{verdeclaro}{\parbox{\dimexpr\linewidth-8pt}{\tiny\strut
\textcolor{verde}{\textbf{\textquestiondown Por qu\'e funciona?}}\\
El hub tiene capacidad para 120 pallets simult\'aneos. Las 12 fincas G20 usan $\sim$70\% del espacio.
El 30\% restante puede servir terceros \textbf{sin CAPEX adicional}. Caso cl\'asico de utilizaci\'on ociosa.\strut}}
\end{column}
\end{columns}

\vspace{4pt}
\colorbox{verdeclaro}{\parbox{\dimexpr\textwidth-12pt}{\tiny\strut
\textcolor{verde}{\textbf{Lectura pedag\'ogica:}} \$652M es \textbf{ingreso} (no ahorro). Cubre el OPEX completo (\$437M) y deja \$215M netos solo de terceros.
Combinado con palancas + transporte: el proyecto se vuelve viable. Fuente: datos Uniban (402 fincas zona), costo de consolidaci\'on \$42 (reuniones).\strut}}
\end{frame}

% ============================================================
% 13f. NUEVAS PALANCAS - EL CASO VIABLE
% ============================================================
\begin{frame}
\frametitle{El hub no solo se paga solo: ahora genera valor}
\tagline{verde}{Nuevas Palancas (Datos 2025)}
\vspace{-2pt}

\begin{center}
\scriptsize
\renewcommand{\arraystretch}{1.2}
\begin{tabular}{c c c c c}
    \toprule
    \textbf{Solo Transporte} & \textbf{+ Standby} & \textbf{+ Nocturnos} & \textbf{=} & \textbf{Neto Palancas} \\
    \midrule
    \textcolor{rojo}{\textbf{$-$\$77M}} & \textcolor{cian}{\textbf{+\$42M}} & \textcolor{cian}{\textbf{+\$100M}} & $=$ & \textcolor{amarillo}{\textbf{\$65M}} \\
    {\tiny Deficit} & {\tiny \$20K/h, 50\%} & {\tiny 50K $\times$ \$4,200} & & {\tiny Insuficiente} \\
    \bottomrule
\end{tabular}
\end{center}

\vspace{6pt}
\begin{columns}[T]
\begin{column}{0.32\textwidth}
\colorbox{amarilloclaro}{\parbox{\dimexpr\linewidth-8pt}{\centering\strut
\textcolor{amarillo}{\textbf{+ Descuento \$75K}}\\[2pt]
{\large\textcolor{amarillo}{\textbf{\$296M}}}\\
{\tiny neto anual}\\[2pt]
Payback: \textcolor{amarillo}{\textbf{13.6 a\~nos}}\strut}}
\end{column}
\begin{column}{0.32\textwidth}
\colorbox{verdeclaro}{\parbox{\dimexpr\linewidth-8pt}{\centering\strut
\textcolor{verde}{\textbf{+ Terceros (10\%)}}\\[2pt]
{\large\textcolor{verde}{\textbf{\$717M}}}\\
{\tiny neto anual}\\[2pt]
Payback: \textcolor{verde}{\textbf{5.6 a\~nos}}\strut}}
\end{column}
\begin{column}{0.32\textwidth}
\colorbox{verdeclaro}{\parbox{\dimexpr\linewidth-8pt}{\centering\strut
\textcolor{verde}{\textbf{Escenario completo}}\\[2pt]
{\large\textcolor{verde}{\textbf{\$948M}}}\\
{\tiny neto anual}\\[2pt]
Payback: \textcolor{verde}{\textbf{4.2 a\~nos}}\strut}}
\end{column}
\end{columns}

\vspace{6pt}
\colorbox{verdeclaro}{\parbox{\dimexpr\textwidth-12pt}{\scriptsize\strut
\textcolor{verde}{\textbf{Conclusi\'on:}} Con CAPEX real (\$4,020M), el proyecto requiere la \textbf{estrategia completa}:
palancas operativas + servicio a terceros. Con 10\% de penetraci\'on regional: payback \textbf{5.6 a\~nos}.
Con todo: \textbf{4.2 a\~nos}.\strut}}
\end{frame}

% ============================================================
% 14. SENSIBILIDAD
% ============================================================
\begin{frame}
\frametitle{\textquestiondown Qu\'e pasa si G20 negocia descuento?}
\tagline{verde}{Sensibilidad (con palancas \$142M incluidas)}
\vspace{-2pt}

{\scriptsize G20 puede negociar hasta \$150K menos por viaje. Tabla incluye palancas operativas (standby + nocturnos):}

\vspace{4pt}
\begin{center}
\scriptsize
\renewcommand{\arraystretch}{1.1}
\begin{tabular}{l r r r r r}
    \toprule
    \textbf{Desc.\ G20} & \textbf{Mula} & \textbf{Hub} & \textbf{Ahorro} & \textbf{Neto/a\~no} & \textbf{Payback} \\
    \midrule
    \rowcolor{rojoclaro}
    \textcolor{rojo}{\textbf{\$0 $\leftarrow$ BASE}} & \$720K & \textbf{\$810} & \$109 & \textcolor{amarillo}{\$65M} & N/A ($>$60a) \\
    \$50K & \$660K & \$763 & \$156 & \textcolor{amarillo}{\$216M} & 18.6a \\
    \$75K & \$630K & \$739 & \$180 & \textcolor{amarillo}{\$296M} & 13.6a \\
    \$100K & \$600K & \$715 & \$204 & \textcolor{verde}{\$376M} & 10.7a \\
    \rowcolor{verdeclaro}
    \textcolor{verde}{\textbf{\$150K max}} & \$540K & \textcolor{verde}{\textbf{\$668}} & \textcolor{verde}{\textbf{\$251}} & \textcolor{verde}{\textbf{\$532M}} & \textcolor{verde}{\textbf{7.6a}} \\
    \bottomrule
\end{tabular}
\end{center}

\vspace{4pt}
\colorbox{verdeclaro}{\parbox{\dimexpr\textwidth-12pt}{\scriptsize\strut
\textcolor{verde}{\textbf{Lectura:}} Con CAPEX real (\$4,020M), el descuento solo no es suficiente.
Pero combinado con palancas (\$142M) y \textbf{terceros} (\$652M), incluso \$0 descuento da
payback \textbf{5.6 a\~nos}. El descuento \textbf{acelera} el retorno.\strut}}

\vspace{1pt}
{\tiny\color{gris} CAPEX \$4,020M. OPEX \$437M/a\~no. 10 mulas/d\'ia $\times$ 260 d\'ias. Palancas: standby \$42M + nocturnos \$100M.}
\end{frame}

% ============================================================
% 14b. SENSIBILIDAD FINCA→HUB
% ============================================================
\begin{frame}
\frametitle{Sensibilidad: costo Finca $\rightarrow$ Hub}
\tagline{amarillo}{Segunda Variable Clave}
\vspace{-2pt}

{\scriptsize El \'acido asume F$\rightarrow$H = \$200/caja, pero el rango real es \$80--\$300. Hub total = \$568 (flete) + \$42 (consol) + F$\rightarrow$H:}

\vspace{3pt}
\begin{center}
\scriptsize
\renewcommand{\arraystretch}{1.05}
\begin{tabular}{l r r r r}
    \toprule
    \textbf{F$\rightarrow$H \$/caja} & \textbf{Hub total} & \textbf{Ahorro} & \textbf{Neto/a\~no} & \textbf{Payback} \\
    \midrule
    \rowcolor{verdeclaro}
    \textcolor{verde}{\textbf{\$80 (optimista)}} & \$690 & \$229 & \textcolor{verde}{\textbf{\$319M}} & \textcolor{verde}{\textbf{12.6a}} \\
    \$120 & \$730 & \$189 & \$187M & 21.5a \\
    \$150 & \$752 & \$167 & \$112M & 35.9a \\
    \rowcolor{rojoclaro}
    \textcolor{rojo}{\textbf{\$200 $\leftarrow$ \'ACIDO}} & \textbf{\$810} & \$109 & \textcolor{rojo}{\textbf{$-$\$77M}} & N/A \\
    \$250 & \$860 & \$59 & $-$\$244M & N/A \\
    \rowcolor{rojoclaro}
    \textcolor{rojo}{\textbf{\$300 (extremo)}} & \$910 & \$9 & $-$\$408M & N/A \\
    \bottomrule
\end{tabular}
\end{center}

\vspace{1pt}
{\tiny\color{gris} Hub total = \$568 + \$42 + F$\rightarrow$H. Ahorro = \$919 $-$ Hub total.
Bruto = ahorro $\times$ 3.295M cajas/a\~no. Neto = bruto $-$ \$437M OPEX.
Payback = \$4,020M $\div$ neto. Solo transporte (sin palancas ni terceros).}

\vspace{4pt}
\colorbox{amarilloclaro}{\parbox{\dimexpr\textwidth-12pt}{\scriptsize\strut
\textcolor{amarillo}{\textbf{Nota:}} Esta sensibilidad es \textbf{con \$0 descuento G20} (\'acido).
Las dos sensibilidades (descuento G20 y F$\rightarrow$H) son \textbf{variables independientes}.\strut}}

\vspace{3pt}
\colorbox{verdeclaro}{\parbox{\dimexpr\textwidth-12pt}{\scriptsize\strut
\textcolor{verde}{\faIcon{sliders-h}}\hspace{2pt}
\textcolor{verde}{\textbf{F$\rightarrow$H es ajustable con el slider en el gemelo digital.}}
Variable clave: negociar tarifa interna corta.\strut}}
\end{frame}

% ============================================================
% 15. CAPEX + OPEX
% ============================================================
\begin{frame}
\frametitle{\textquestiondown Cu\'anto cuesta construir el centro?}
\tagline{azul}{Inversi\'on real (Costos 2026)}
\vspace{-2pt}

\begin{columns}[T]
\begin{column}{0.48\textwidth}
    {\tiny\color{gris}\textbf{CAPEX --- INFRAESTRUCTURA REAL (1,800 m$^2$)}}
    \vspace{2pt}

    \tiny
    \renewcommand{\arraystretch}{1.05}
    \begin{tabular}{l r}
        \toprule
        Lote 5,000 m$^2$ (Chinita) & \$200M \\
        Bodega consolidaci\'on (1,000 m$^2$) & \$1,900M \\
        Plataformas + 4 muelles (250 m$^2$) & \$475M \\
        P\'ergolas + circulaciones (400 m$^2$) & \$640M \\
        Oficinas + servicios (150 m$^2$) & \$405M \\
        Equipos (dock levelers, montacargas) & \$250M \\
        Capital de trabajo & \$150M \\
        \midrule
        \textcolor{azul}{\textbf{TOTAL CAPEX}} & \textcolor{azul}{\textbf{\$4,020M}} \\
        \bottomrule
    \end{tabular}

    \vspace{3pt}
    {\tiny\color{gris} Bodega: \$1.6--\$2.2M/m$^2$. Admin: \$2.2--\$3.0M/m$^2$.
    Incluye indirectos 15--20\% (dise\~no, licencias, interventor\'ia, imprevistos).}
\end{column}

\begin{column}{0.48\textwidth}
    {\tiny\color{gris}\textbf{OPEX ANUAL (escalado)}}
    \vspace{2pt}

    \tiny
    \renewcommand{\arraystretch}{1.05}
    \begin{tabular}{l r}
        \toprule
        Personal (6) & \$215M \\
        Servicios & \$73M \\
        Mantenimiento (2\% construcci\'on) & \$73M \\
        Seguros (1\% CAPEX) & \$40M \\
        Administraci\'on & \$36M \\
        \midrule
        \textcolor{amarillo}{\textbf{TOTAL OPEX}} & \textcolor{amarillo}{\textbf{\$437M/a\~no}} \\
        \bottomrule
    \end{tabular}

    \vspace{4pt}
    \colorbox{azulclaro}{\parbox{\dimexpr\linewidth-8pt}{\tiny\strut
    \textcolor{azul}{\textbf{Especificaciones:}}\\
    4 muelles (2 recepci\'on + 2 despacho 24p)\\
    Patio maniobra 1,500 m$^2$, radio giro 12 m\\
    Bodega altura 6--8 m, dock levelers\\
    Preinstalaci\'on cadena fr\'io\strut}}
\end{column}
\end{columns}

\vspace{3pt}
{\tiny\color{gris} Fuente: INPUTS CENTRO DE CONSOLIDACI\'ON (Arquitectura, Abr 2026). Anterior estimado: \$1,200M.}
\end{frame}

% ============================================================
% 16. BENEFICIOS ESTRATÉGICOS
% ============================================================
\begin{frame}
\frametitle{M\'as que ahorro en transporte}
\tagline{azul}{Beneficios Estrat\'egicos}
\vspace{-2pt}

\begin{columns}[T]
\begin{column}{0.30\textwidth}
\colorbox{verdeclaro}{\parbox{\dimexpr\linewidth-8pt}{\scriptsize\strut
\textcolor{verde}{\textbf{\Large $-$39\%}}\\[2pt]
\textbf{V\'ia principal}\\[3pt]
41 $\rightarrow$ 25 en v\'ia ppal.\\
+26 cortos finca$\rightarrow$hub.\strut}}
\end{column}
\begin{column}{0.30\textwidth}
\colorbox{azulclaro}{\parbox{\dimexpr\linewidth-8pt}{\scriptsize\strut
\textcolor{azul}{\textbf{\Large $<$1h}}\\[2pt]
\textbf{Espera en puerto}\\[3pt]
Contenerizado listo\\
vs 5--6h hoy.\strut}}
\end{column}
\begin{column}{0.30\textwidth}
\colorbox{amarilloclaro}{\parbox{\dimexpr\linewidth-8pt}{\scriptsize\strut
\textcolor{amarillo}{\textbf{\Large 402}}\\[2pt]
\textbf{Fincas en zona}\\[3pt]
Servicio a terceros.\\
Ingresos adicionales.\strut}}
\end{column}
\end{columns}

\vspace{4pt}
\begin{columns}[T]
\begin{column}{0.47\textwidth}
    {\scriptsize\color{rojo}\textbf{HOY (SIN CENTRO)}}
    \begin{itemize}\scriptsize
        \setlength\itemsep{-2pt}
        \item[\textcolor{rojo}{\scriptsize$\blacksquare$}] Recolecci\'on a las 2--3~AM
        \item[\textcolor{rojo}{\scriptsize$\blacksquare$}] Camiones devueltos a las 11~PM
        \item[\textcolor{rojo}{\scriptsize$\blacksquare$}] Coordinaci\'on manual 19 fincas
        \item[\textcolor{rojo}{\scriptsize$\blacksquare$}] Sin trazabilidad
    \end{itemize}
\end{column}
\begin{column}{0.47\textwidth}
    {\scriptsize\color{verde}\textbf{CON CENTRO}}
    \begin{itemize}\scriptsize
        \setlength\itemsep{-2pt}
        \item Turnos planificados desde 9~AM
        \item Fruta contenerizada
        \item Trazabilidad pallet $\rightarrow$ puerto
        \item Escalable a terceros
    \end{itemize}
\end{column}
\end{columns}
\end{frame}

% ============================================================
% 16b. DATOS GPS / GEOCERCAS
% ============================================================
\begin{frame}
\frametitle{Datos GPS reales: 9--12 Marzo 2026}
\tagline{azul}{Evidencia Emp\'irica}
\vspace{-2pt}

{\scriptsize Geocercas instaladas en 7 veh\'iculos del Grupo 20. Datos brutos que alimentan el modelo:}

\vspace{4pt}
\begin{columns}[T]
\begin{column}{0.47\textwidth}
{\scriptsize\color{azul}\textbf{RESUMEN DE RASTREO}}
\vspace{2pt}

\scriptsize
\renewcommand{\arraystretch}{1.1}
\begin{tabular}{>{\color{azul}\bfseries}r l}
    \toprule
    7 & veh\'iculos rastreados \\
    42 & viajes registrados \\
    2.5 & viajes/d\'ia/veh\'iculo \\
    78 min & promedio en embarcadero \\
    45 min & promedio cargue en finca \\
    22 min & Casa Verde $\rightarrow$ Zungo \\
    68 km/h & velocidad m\'axima en v\'ia \\
    \bottomrule
\end{tabular}
\end{column}

\begin{column}{0.47\textwidth}
{\scriptsize\color{verde}\textbf{UBICACIONES CLAVE}}
\vspace{2pt}

\scriptsize
\renewcommand{\arraystretch}{1.05}
\begin{tabular}{l l}
    \toprule
    \textbf{Punto} & \textbf{Rol} \\
    \midrule
    Casa Verde & Base operativa \\
    Zungo & Embarcadero fluvial \\
    Chinita & Ubicaci\'on hub \\
    Puerto Antioqu. & Destino final \\
    Fincas (19) & Origen de carga \\
    \bottomrule
\end{tabular}

\vspace{4pt}
\colorbox{grisclaro}{\parbox{\dimexpr\linewidth-8pt}{\tiny\strut
\textcolor{gris}{\textbf{Fuente:}} IP Uniban, GPS\\
Anderson Mestra, Marzo 2026\strut}}
\end{column}
\end{columns}

\vspace{4pt}
\colorbox{verdeclaro}{\parbox{\dimexpr\textwidth-12pt}{\scriptsize\strut
\textcolor{verde}{\textbf{Estos datos calibran el gemelo digital:}} tiempos de cargue, velocidades,
esperas en embarcadero y frecuencia de viajes. El modelo no asume --- \textbf{mide}.\strut}}
\end{frame}

% ============================================================
% 17. ESCALABILIDAD (piso corregido)
% ============================================================
\begin{frame}
\frametitle{Servicio a terceros: acelerando el payback}
\tagline{cian}{Escalabilidad}
\vspace{-2pt}

{\scriptsize G20 = 4.7\% regi\'on. Consolidaci\'on cuesta \$42/caja $\rightarrow$ tarifa m\'inima a terceros: \textbf{$>$\$42}.}

\vspace{4pt}
\begin{center}
\scriptsize
\renewcommand{\arraystretch}{1.1}
\begin{tabular}{r r r r}
    \toprule
    & \multicolumn{3}{c}{\textbf{Ingreso anual por servicio a terceros}} \\
    \cmidrule(lr){2-4}
    \textbf{Tarifa/caja} & \textbf{5\% regi\'on} & \textbf{10\% regi\'on} & \textbf{25\% regi\'on} \\
    \midrule
    \rowcolor{rojoclaro}
    {\color{gris}\$42 (= costo)} & {\color{gris}\$228M} & {\color{gris}\$457M} & {\color{gris}\$1,142M} \\
    \$50  & \$272M & \$544M & \$1,359M \\
    \$60  & \$326M & \textcolor{verde}{\textbf{\$652M}} & \$1,631M \\
    \$80  & \$435M & \$870M & \$2,175M \\
    \bottomrule
\end{tabular}
\end{center}

\vspace{4pt}
\colorbox{verdeclaro}{\parbox{\dimexpr\textwidth-12pt}{\scriptsize\strut
\textcolor{verde}{\textbf{Con \$60/caja y 10\% regi\'on:}} ingreso \$652M \textbf{$>$ OPEX \$437M}.
Combinado con transporte + palancas: neto \textbf{\$717M/a\~no}, payback \textbf{5.6 a\~nos}.
\textbf{Terceros son la pieza clave de viabilidad} con CAPEX real (\$4,020M).\strut}}

\vspace{1pt}
{\tiny\color{gris} PA 7M ton. Uniban 250+ fincas. 402 bananeras en zona. Piso = costo consolidaci\'on \$42/caja.}
\end{frame}

% ============================================================
% 18. GEMELO DIGITAL
% ============================================================
\begin{frame}
\frametitle{Gemelo Digital Interactivo}
\tagline{verde}{La Herramienta}

\vspace{6pt}
\begin{center}
\begin{tikzpicture}
    \node[rounded corners=6pt, fill=grisclaro, draw=gris!30,
          minimum width=10cm, minimum height=2.8cm] (s) {};
    \node at (s.center) {
        \begin{tabular}{c}
            \textcolor{verde}{\Large\faGlobeAmericas}\\[3pt]
            {\large\textbf{gemelo-digital-uraba.vercel.app}}\\[3pt]
            {\footnotesize\color{gris} SC1 Hoy $\cdot$ SC3 Con Hub $\cdot$ Sensibilidad G20}\\[2pt]
            {\footnotesize\color{gris} Click por finca $\cdot$ Tabla de sensibilidad $\cdot$ Dashboard financiero}
        \end{tabular}
    };
\end{tikzpicture}
\end{center}

\vspace{6pt}
\begin{center}
\colorbox{verde}{\small\bfseries\color{white}\strut\hspace{8pt}ABRIR GEMELO DIGITAL\hspace{8pt}}
\end{center}

\vspace{4pt}
{\scriptsize\color{gris} Todo lo presentado est\'a vivo en la herramienta. Se pueden ajustar par\'ametros (producci\'on, tarifas, finca$\rightarrow$hub) y ver resultados en tiempo real. 19 fincas, 39 rutas OSRM, geocercas GPS.}
\end{frame}

% ============================================================
% 19. SIGUIENTES PASOS
% ============================================================
\begin{frame}
\frametitle{Siguientes pasos}
\tagline{verde}{Para la Junta}
\vspace{-2pt}

\begin{block}{\scriptsize Resumen con CAPEX real \$4,020M, OPEX \$437M/a\~no}
\scriptsize
\textbf{Sin hub:} \$919/caja, 26 sencillos Centro, 41 viajes/d\'ia.\\
\textbf{Con hub (\'acido, solo transporte):} \$810/caja, 10 mulas. Neto \textbf{$-$\$77M/a\~no} (deficit).\\
\textbf{+ Palancas operativas (\$142M):} standby \$42M + nocturnos \$100M $\rightarrow$ neto \textbf{\$65M/a\~no}.\\
\textbf{+ Terceros (10\%, \$60/caja, \$652M):} payback \textbf{5.6 a\~nos}.\\
\textbf{+ Descuento \$75K mula + terceros + palancas:} payback \textbf{4.2 a\~nos}, neto \$948M/a\~no.\\
\textbf{V\'ia ppal:} $-$39\% (41$\rightarrow$25), +26 cortos finca$\rightarrow$hub. Contenerizado, $<$1h espera.
\end{block}

\vspace{4pt}
\begin{center}
\begin{tikzpicture}
    \node[rounded corners=4pt, fill=verdeclaro, draw=verde!30,
          minimum width=2.3cm, minimum height=0.9cm] (s1) {};
    \node[font=\scriptsize\bfseries, text=verde] at ([yshift=2pt]s1.center) {1. CAPEX};
    \node[font=\tiny, text=gris, below=0pt of s1.center] {Cotizar obra};

    \node[rounded corners=4pt, fill=azulclaro, draw=azul!30,
          minimum width=2.3cm, minimum height=0.9cm, right=4pt of s1] (s2) {};
    \node[font=\scriptsize\bfseries, text=azul] at ([yshift=2pt]s2.center) {2. COSTEO};
    \node[font=\tiny, text=gris, below=0pt of s2.center] {Costos por finca};

    \node[rounded corners=4pt, fill=amarilloclaro, draw=amarillo!30,
          minimum width=2.3cm, minimum height=0.9cm, right=4pt of s2] (s3) {};
    \node[font=\scriptsize\bfseries, text=amarillo] at ([yshift=2pt]s3.center) {3. NEGOCIAR};
    \node[font=\tiny, text=gris, below=0pt of s3.center] {Desc.\ mula G20};

    \node[rounded corners=4pt, fill=grisclaro, draw=gris!30,
          minimum width=2.3cm, minimum height=0.9cm, right=4pt of s3] (s4) {};
    \node[font=\scriptsize\bfseries, text=gris] at ([yshift=2pt]s4.center) {4. TIMELINE};
    \node[font=\tiny, text=gris, below=0pt of s4.center] {Cronograma};
\end{tikzpicture}
\end{center}

\vspace{6pt}
\begin{center}
\colorbox{verde}{\normalsize\bfseries\color{white}\strut\hspace{8pt}VER GEMELO DIGITAL EN VIVO\hspace{8pt}}
\vspace{2pt}

{\tiny\color{gris} Observatorio de Datos y An\'alisis SAS $\cdot$ Anderson Mestra, Grupo 20 $\cdot$ Marzo 2026}
\end{center}
\end{frame}

\end{document}
